% ===== グラフ設定ファイル =====
% グラフに関するすべての設定とマクロをここに記述
% main.tex から \input{figure_config.tex} で読み込まれます

% ---- pgfplots/tikz 関連パッケージ ----
\usepackage{pgfplots}
\usepackage{pgfplotstable}
\usepackage{tikz}
\usepgfplotslibrary{fillbetween}

% ---- グラフマーカーの定義 ----
\pgfplotscreateplotcyclelist{reportcyclelist}{%
    {black, mark=o, mark size=3.2pt, line width=0.9pt, mark options={draw=black, fill=none}},%
    {black, mark=square, mark size=3.2pt, line width=0.9pt, mark options={draw=black, fill=none}},%
    {black, mark=triangle, mark size=3.2pt, line width=0.9pt, mark options={draw=black, fill=none}},%
    {black, mark=diamond, mark size=3.2pt, line width=0.9pt, mark options={draw=black, fill=none}},%
    {black, mark=x, mark size=3.2pt, line width=0.9pt}%
}

% ---- グラフ全体の設定 ----
\pgfplotsset{
    compat=1.18,
    every axis/.style={
        width=10cm,
        height=10cm,
        axis lines=box,
        axis line style={black, line width=1.0pt},
        tick style={black, line width=0.8pt},
        tick align=inside,
        tick label style={font=\small},
        label style={font=\small},
        title style={font=\small\bfseries},
        grid=none,
        minor tick num=1,
        scaled ticks=false,
        legend cell align=left,
        legend style={
            draw=black,
            fill=white,
            fill opacity=0.85,
            text opacity=1,
            font=\small,
            rounded corners=1pt,
        },
        cycle list name=reportcyclelist,
    },
    plot area/.style={width=10cm, height=10cm},
    plot square/.style={width=10cm, height=10cm},
    plot logx/.style={xmode=log},
    plot logy/.style={ymode=log},
    plot loglog/.style={xmode=log, ymode=log},
    plot no grid/.style={grid=none},
    fit line/.style={black, no marks, line width=1.1pt, on layer=axis foreground},
    scatter points/.style={only marks, mark=o, mark size=3.2pt, mark options={draw=black, fill=none}, on layer=axis background},
    line series/.style={line width=1.0pt},
    area series/.style={fill opacity=0.18, on layer=axis background},
}

% ========================================
% 補助マクロ
% ========================================

\newcommand{\PlotTable}[4][]{%
    \addplot[scatter points, #1] table [x=#2, y=#3, col sep=comma] {#4};%
}

\newcommand{\PlotFitLine}[4]{%
    % #1=slope, #2=intercept, #3=xmin, #4=xmax
    \addplot[fit line] coordinates {(#3,{#1*(#3)+#2}) (#4,{#1*(#4)+#2})};%
}

% ========================================
% グラフ定義（tikzpicture で描画）
% 以下、各グラフをマクロ化して定義
% ========================================

% ---- 図1: 基本プロット ----
\newcommand{\GraphBasicPlot}{%
    \begin{tikzpicture}
    \begin{axis}[
        plot area,
        title={電圧測定結果},
        xlabel={時間 [s]},
        ylabel={電圧 [V]},
        legend style={at={(0.98,0.98)}, anchor=north east},
    ]
        \addplot[scatter points] coordinates {
            (0, 0)
            (1, 1.2)
            (2, 2.1)
            (3, 3.0)
            (4, 3.8)
            (5, 4.9)
        };
        \addlegendentry{測定値}

        \addplot[fit line] coordinates {
            (0, 0)
            (5, 5.1)
        };
        \addlegendentry{近似直線}
    \end{axis}
    \end{tikzpicture}
}

% ---- 図2: 温度特性 ----
\newcommand{\GraphTemperatureResponse}{%
    \begin{tikzpicture}
    \begin{axis}[
        plot area,
        title={抵抗の温度特性},
        xlabel={温度 [$^\circ$C]},
        ylabel={抵抗 [$\Omega$]},
        legend style={at={(0.98,0.02)}, anchor=south east},
    ]
        \addplot[scatter points] table [x=temperature, y=resistance, col sep=comma] {data/temperature_response.csv};
        \addlegendentry{実験値}

        \addplot[fit line] coordinates {
            (0, 1000)
            (100, 500)
        };
        \addlegendentry{近似直線}
    \end{axis}
    \end{tikzpicture}
}

% ---- 図3: 周波数応答 ----
\newcommand{\GraphFrequencyResponse}{%
    \begin{tikzpicture}
    \begin{axis}[
        plot area,
        title={周波数応答（ゲイン）},
        xlabel={周波数 [Hz]},
        ylabel={ゲイン},
        xmode=log,
        legend style={at={(0.98,0.02)}, anchor=south east},
    ]
        \addplot[scatter points] table [x=frequency, y=magnitude, col sep=comma] {data/frequency_response.csv};
        \addlegendentry{測定値}

        \addplot[fit line] coordinates {
            (1, 0.95)
            (10000, 0.15)
        };
        \addlegendentry{近似曲線}
    \end{axis}
    \end{tikzpicture}
}
